%% kapa: exam=1 round=36 year=2025 period=1 form=A subject=real_estate subject_ko=부동산학원론 seq=10 number=90
%% source: Q-Net 감정평가사 기출문제 (cst00309) artlSeq=5231525 file=raw/1008/5231525/2025년 제36회 감정평가사 1차시험 1교시 문제지 원본.pdf
%% conversion: kordoc 4.24.51 -> tools/gen1.py (auto); review=proofread-sonnet; flags=none
\begin{question}{90}
A지역 아파트시장의 기존 수요함수는 $2P=-Qd+400$ , 공급함수는 $P_{1}=Qs_{1}+20$ 이었다. 이후 수요함수는 변하지 않고 공급함수가 $P_{2}=Qs_{2}+80$ 으로 변하였다. 다음 설명으로 옳은 것은? (단, $X$ 축은 수량, $Y$ 축은 가격, $P$ 는 가격 (단위: 만원/ $\mathrm{m}^{2}$ ), $Qd$ 는 수요량(단위: $\mathrm{m}^{2}$ ), $Qs$ 는 공급량(단위: $\mathrm{m}^{2}$ )이며, 다른 조건은 동일함)
\begin{choices}
  \item 아파트 공급량의 증가에 따라 공급곡선이 좌측(좌상향)으로 이동한다.
  \item 기존 아파트시장의 균형가격은 120만원/ $\mathrm{m}^{2}$이다.
  \item 공급함수 변화이후, 아파트시장의 균형거래량은 160 $\mathrm{m}^{2}$이다.
  \item 공급함수 변화이후, 아파트시장의 균형가격은 20만원/ $\mathrm{m}^{2}$ 만큼 감소한다.
  \item 공급함수 변화이후, 아파트시장의 균형거래량은 40 $\mathrm{m}^{2}$ 만큼 감소한다.
\end{choices}
\end{question}
\answer{5}
